\documentclass[11pt]{article}
\usepackage[margin=1in]{geometry}
\usepackage{amsmath,amssymb}
\usepackage{graphicx}
\usepackage{hyperref}
\usepackage{booktabs}

\title{Notebook 11: Learned Gap Model}
\author{prime-numbers-lab}
\date{}

\begin{document}
\maketitle

\section*{Purpose}
Notebook 11 tests whether replacing the baseline expected prime gap model
\[
  g(x) \approx \log x
\]
with learned deterministic alternatives materially improves reconstruction diagnostics.

\section*{Models}
The notebook compares five expected-gap models:
\begin{itemize}
  \item log baseline: $g(x)=\log x$
  \item affine log: $g(x)=a\log x+b$
  \item sinusoidal log correction: $g(x)=\log x + A\sin(k\log x)+B\cos(k\log x)+c$
  \item local rolling residual correction
  \item combined sinusoidal plus local residual correction
\end{itemize}

\section*{Primary diagnostics}
The primary diagnostics are global RMSE/MAE and windowed residual drift.
The best model is \texttt{local_rolling} with RMSE 10.224070.
The log-baseline RMSE is 10.225042.
The best RMSE improvement versus log baseline is 0.0095%.

\section*{Figures}
\begin{itemize}
  \item \texttt{figures/11\_gap\_model\_rmse.png}
  \item \texttt{figures/11\_gap\_model\_mae.png}
  \item \texttt{figures/11\_expected\_gap\_curves.png}
  \item \texttt{figures/11\_window\_relative\_residual\_drift.png}
  \item \texttt{figures/11\_window\_mean\_residual.png}
  \item \texttt{figures/11\_rmse\_improvement\_vs\_log.png}
\end{itemize}

\section*{Interpretation}
At this scale, deterministic expected-gap refinements produce only tiny improvements over
$\log x$. This indicates that residual gap error is dominated more by distributional
variance than by a smooth mean-model bias.

\section*{Next step}
Notebook 12 should model
\[
  P(g\mid x)
\]
rather than only a point estimate $\hat g(x)$.

\section*{Constraint}
\[
\text{Constraint} \rightarrow \text{signal} > \text{noise}
\]

\end{document}
